%%%%%%%%%%%%%%%%%%%%%%%%%%%%%%%%%%%%%%%%%%%%%%%%%%%%%%%%%%%%%%%%%%%%%%
% Lớp 10 | Véctơ (chưa xét tọa độ) | Tích vô hướng (chưa xét tọa độ)
% Dạng 0H5?4-1: Tích vô hướng, góc giữa 2 véctơ
% Số câu: 36  (sinh tự động bởi tach_cau.py; từ nay thầy sửa trực tiếp file này)
%%%%%%%%%%%%%%%%%%%%%%%%%%%%%%%%%%%%%%%%%%%%%%%%%%%%%%%%%%%%%%%%%%%%%%

%%%%%%%%%%%%%%%% TRẮC NGHIỆM 4 PHƯƠNG ÁN %%%%%%%%%%%%%%%%

% nguồn: MaHoaNH_TracNghiem_Lop10
\begin{ex}%[0H5N4-1]%[ma:b5da166b]
	Trong các khẳng định sau, khẳng định nào đúng với mọi $\overrightarrow a$ và $\overrightarrow  b$?
	 \choice
		{\True $\left| \overrightarrow a\cdot \overrightarrow b  \right|=\left| \overrightarrow a \right|\left| \overrightarrow b \right|\left| \cos\left( \overrightarrow a,\overrightarrow b \right) \right|$}
		{$\overrightarrow a\cdot \overrightarrow b=-\left| \overrightarrow  a \right|\left| \overrightarrow b \right|$}
		{$\overrightarrow a\cdot \overrightarrow b=\left| \overrightarrow a \right|\left| \overrightarrow b \right|$} 
		{$\left| \overrightarrow a\cdot \overrightarrow b  \right|=\left| \overrightarrow a \right|\left| \overrightarrow b \right|$}
	 \loigiai{
		 Theo định nghĩa tích vô hướng hai vectơ, ta có \[\overrightarrow a\cdot \overrightarrow b=\left| \overrightarrow a \right|\left| \overrightarrow b \right|\cos\left( \overrightarrow a,\overrightarrow b \right)\Rightarrow \left| \overrightarrow a \cdot \overrightarrow b \right|=\left| \overrightarrow a \right|\left| \overrightarrow b \right|\left| \cos\left( \overrightarrow a,\overrightarrow b \right) \right|.\]
	 }
\end{ex}

% nguồn: MaHoaNH_TracNghiem_Lop10
\begin{ex}%[0H5N4-1]%[ma:59567f68]
	Cho hình vuông $ABCD$ có độ dài cạnh bằng $2$.  Tính $\overrightarrow{AB} \cdot \overrightarrow{AD}$.
	 \choice
		 {\True $\overrightarrow{AB} \cdot \overrightarrow{AD}=0$}
		{$\overrightarrow{AB} \cdot \overrightarrow{AD}=4$} 
		 {$\overrightarrow{AB} \cdot \overrightarrow{AD}=\overrightarrow 0$}
		{$\overrightarrow{AB} \cdot \overrightarrow{AD}=-4$} 
	 \loigiai{
		Vì $ABCD$ là hình vuông nên $AB\perp AD$. Do đó, $\overrightarrow{AB} \cdot  \overrightarrow{AD}=0$.
	 }
\end{ex}

% nguồn: MaHoaNH_TracNghiem_Lop10
\begin{ex}%[0H5N4-1]%[ma:dd119df6]
	Cho tam giác đều $ABC$ có cạnh bằng  $a$. Tính $\overrightarrow{AB} \cdot \overrightarrow{AC}$.
	 \choice
		{$\overrightarrow{AB} \cdot \overrightarrow{AC}=\dfrac{\sqrt 3}{2}a^2$} 
		{\True $\overrightarrow{AB} \cdot \overrightarrow{AC}=\dfrac{1}{2}a^2$} 
		{$\overrightarrow{AB} \cdot \overrightarrow{AC}=-\dfrac{1}{2}a^2$} 
		{$\overrightarrow{AB} \cdot  \overrightarrow{AC}=a^2$}
	\loigiai{ 
		Ta có $\overrightarrow{AB} \cdot \overrightarrow{AC}$ $=AB \cdot AC \cdot \cos \widehat{BAC}$ $=a \cdot a  \cdot \cos 60^\circ $ $=\dfrac{1}{2}a^2$.
	} 
\end{ex}

% nguồn: MaHoaNH_TracNghiem_Lop10
\begin{ex}%[0H5N4-1]%[ma:fa09f28a]
	Cho tam giác $\Delta ABC$ đều cạnh $a$. Tích  vô hướng $\overrightarrow{AB} \cdot \overrightarrow{BC}$ bằng
	\choice 
		{$a^2$} 
		{$\dfrac{a^2\sqrt  3} 2$}
		 {\True $-\dfrac{a^2}{2}$}
		{$\dfrac{a^2}{2}$} 
	\loigiai{ 
	 \immini{
		Dựng $\overrightarrow{BD}=\overrightarrow{AB}$.\\ 
		Ta có $\left(  \overrightarrow{AB},\overrightarrow{BC} \right)=\left( \overrightarrow{BC},\overrightarrow{BD} \right)=\widehat{CBD}=120^\circ$\\
		Vậy $\overrightarrow{AB} \cdot \overrightarrow{BC}=AB \cdot BC\cdot  \cos120^\circ=-\dfrac{a^2}{2}$.}{
		\begin{tikzpicture}[line join = round,line cap = round,scale=1]
				 \def \a{2};
				\coordinate (B)  at (0,0);
				\coordinate (C) at ($(B)+(\a,0)$); 
				\coordinate[shift=(60:\a cm)] (A) at  (B);
				\coordinate (D) at  ($(A)!2!(B)$);
				\draw pic[draw,angle  eccentricity=1.4,pic text=$ $,angle radius=0.35cm]{angle=D--B--C};
				 \draw (A)--(B)--(C)--cycle;
				\draw[->,arrows={-latex}]  (B)--(D);
				\draw[->,arrows={-latex}]  (A)--(B);
				\draw[->,arrows={-latex}]  (B)--(C);
				\foreach  \x/\y in {A/C,A/B,B/C}{
					 \path (\x)--(\y) node[midway,sloped]{\tikz{\draw (-90:2pt)--(90:2pt);}};
					 \draw (\x)--(\y) node[midway,sloped] {};
				} 
				\foreach  \x/\y in {A/90,B/180,C/0,D/-180}
				\draw[fill=black] (\x) circle (0.5pt) +  (\y:0.5cm) node{$\x$};
			\end{tikzpicture}
		} 
	} 
\end{ex}

% nguồn: MaHoaNH_TracNghiem_Lop10
\begin{ex}%[0H5N4-1]%[ma:88149a75]
	Cho tam giác $ABC$ có $\widehat{A}=90^\circ,\widehat{B}=60^\circ$ và $AB=a$. Khi đó $\overrightarrow{AC} \cdot \overrightarrow{CB}$  bằng
	\choice 
		 {$-2a^2$}
		{$2a^2$} 
		 {$3a^2$}
		 {\True $-3a^2$}
	 \loigiai{
	\immini{ 
		 Dựng $\overrightarrow{CA'}=\overrightarrow{AC}$. Suy ra $\widehat{BCA'}=180^\circ-30^\circ=150^\circ$.\\
		Trong tam giác  $ABC$ vuông tại $A$, ta có $$AC=AB\cdot \tan 60^\circ=a\sqrt{3}.$$
		Ta có $\left(  \overrightarrow{AC},\overrightarrow{CB} \right)=\left( \overrightarrow{CA'},\overrightarrow{CB} \right)=\widehat{BCA'}=150^\circ$\\
		Vậy \begin{eqnarray*}
		 \overrightarrow{AC} \cdot \overrightarrow{CB}&=&\overrightarrow{CA'} \cdot \overrightarrow{CB}\\
		 &=&CA\cdot CB\cdot \cos \widehat{BCA'}\\
		&=&\sqrt 3a\cdot 2a\cdot \left(-\dfrac{\sqrt 3}  2\right)\\
		 &=&-3a^2.
	\end{eqnarray*}
	 }{
	\begin{tikzpicture}[scale=0.7,font=\footnotesize,line join=round,line cap=round,>=stealth]
			\path 
				(-2,0) coordinate  (B)
				 (2,0) coordinate (C)
				(120:2) coordinate  (A);
			\coordinate (A')  at ($(C)+(C)-(A)$);
			\draw  (A)--(B)--(C)--(A)
				(C)--(A'); 
			\node  at ($(B)+(0.15,0.35)$) [right]{$60^\circ$};
			 \draw pic[draw,angle eccentricity=0.6,pic text=$ $,angle radius=0.25cm]{angle=C--B--A};
			\draw pic[draw,angle eccentricity=1.2,pic text=$ $,angle  radius=0.25cm]{angle=B--C--A'};
			\node at ($(C)+(-0.85,-0.55)$) [right]{$150^\circ$}; 
			 \draw pic[draw,angle radius = 6pt]{right angle = B--A--C};
			\foreach \p/\q in {A/90,B/-90,C/30,A'/0} \fill[black] (\p) circle  (1pt)
				($(\p)+(\q:2.5mm)$)  node{$\p$};
		\end{tikzpicture}
	} 
 }
\end{ex}

% nguồn: MaHoaNH_TracNghiem_Lop10
\begin{ex}%[0H5H4-1]%[ma:a8fb7e4a]
	Cho tam giác đều $ABC$ có  $AB=a$. Tính tích vô hướng $\overrightarrow{AB}\cdot \overrightarrow{CA}$.
	\choice 
		{$\overrightarrow{AB}\cdot\overrightarrow{CA}=-\dfrac{\sqrt{3}}{2}a^2$} 
		 {$\overrightarrow{AB}\cdot\overrightarrow{CA}=\dfrac{\sqrt{3}}{2}a^2$}
		{\True  $\overrightarrow{AB}\cdot\overrightarrow{CA}=-\dfrac{1}{2}a^2$}
		{$\overrightarrow{AB}\cdot\overrightarrow{CA}=\dfrac{1}{2}a^2$} 
	\loigiai{ 
	Ta có $\overrightarrow{AB}\cdot\overrightarrow{CA}=-\overrightarrow{AB}\cdot\overrightarrow{AC}=-AB\cdot AC\cdot \cos \widehat{BAC}=-a\cdot  a\cdot \dfrac{1}{2}=-\dfrac{1}{2}a^2$.}
\end{ex}

% nguồn: MaHoaNH_TracNghiem_Lop10
\begin{ex}%[0H5H4-1]%[ma:364845ad]
	Cho tam  giác đều $ABC$. Tính $\left(\overrightarrow{BA}, \overrightarrow{CB}\right)$.
	\choice 
		 {$60^{\circ}$}
		{\True  $120^{\circ}$}
		 {$90^{\circ}$}
		{$180^{\circ}$} 
	\loigiai{ 
	Ta có  $\left(\overrightarrow{BA},\overrightarrow{CB}\right)=\left(\overrightarrow{BA};-\overrightarrow{BC}\right)=180^\circ-\left(\overrightarrow{BA},\overrightarrow{BC}\right)=180^{\circ}-\widehat{ABC}=180^\circ-60^\circ=120^\circ$.}
\end{ex}

% nguồn: MaHoaNH_TracNghiem_Lop10
\begin{ex}%[0H5H4-1]%[ma:e7176f59]
	Cho tam giác $ABC$ có $AB=2$ cm, $BC=4$ cm, $CA=5$  cm. Tính $\overrightarrow{CA}\cdot\overrightarrow{CB}$.
	\choice 
		{$\overrightarrow{CA}\cdot\overrightarrow{CB}=37$} 
		{\True $\overrightarrow{CA}\cdot \overrightarrow{CB}=\dfrac{37}{2}$} 
		{ $\overrightarrow{CA}\cdot  \overrightarrow{CB}=\dfrac{37}{20}$}
		 {$\overrightarrow{CA}\cdot \overrightarrow{CB}=-\dfrac{37}{2}$}
	 \loigiai{
		 Ta có $\overrightarrow{CA}\cdot\overrightarrow{CB}=\dfrac{1}{2}\left(CA^2+CB^2-AB^2\right)=\dfrac{1}{2}\left(5^2+4^2-2^2\right)=\dfrac{37}{2}$.
	 }
\end{ex}

% nguồn: MaHoaNH_TracNghiem_Lop10
\begin{ex}%[0H5H4-1]%[ma:67bb4e56]
	Cho tam giác $ABC$ vuông ở  $A$ có góc $\widehat{ABC}=60^{\circ}$. Số đo góc giữa hai vectơ $\overrightarrow{AC}$ và $\overrightarrow{CB}$ bằng
	 \choice
		{$30^{\circ}$} 
		{$60^{\circ}$} 
		{$120^{\circ}$} 
		{\True  $150^{\circ}$}
	\loigiai{ 
		Ta có $\widehat{ACB}=180^\circ-90^\circ-60^\circ=30^\circ$.\\ 
		 Khi đó, $\left(\overrightarrow{AC},\overrightarrow{CB}\right)=\left(-\overrightarrow{CA},\overrightarrow{CB}\right)=180^\circ-\left(\overrightarrow{CA},\overrightarrow{CB}\right)=180^\circ-\widehat{ACB}=150^\circ$.
	} 
\end{ex}

% nguồn: MaHoaNH_TracNghiem_Lop10
\begin{ex}%[0H5H4-1]%[ma:9ee4bf99]
	Cho tam giác $ABC$ với  $\widehat{A}=60^{\circ}$. Tính tổng $\left(\overrightarrow{AB}, \overrightarrow{BC}\right)+\left(\overrightarrow{BC}, \overrightarrow{CA}\right)$.
	\choice 
		 {$120^{\circ}$}
		{$360^{\circ}$} 
		{$270^{\circ}$} 
		 {\True $240^{\circ}$}
	\loigiai{ 
		Ta có \begin{align*}
		 \left(\overrightarrow{AB}, \overrightarrow{BC}\right)+\left(\overrightarrow{BC}, \overrightarrow{CA}\right)&=\left(-\overrightarrow{BA}, \overrightarrow{BC}\right)+\left(-\overrightarrow{CB}, \overrightarrow{CA}\right)\\
		&=180^\circ  -\left(\overrightarrow{BA}, \overrightarrow{BC}\right)+180^\circ-\left(\overrightarrow{CB}, \overrightarrow{CA}\right)\\
		 &=360^\circ-\left(\widehat{ABC}+\widehat{ACB}\right)\\
		&=360^\circ-\left(180^\circ-\widehat{BAC}\right)\\ 
		&=360^\circ-\left(180^\circ-60^\circ\right)\\ 
		 &=240^\circ.
	\end{align*}
} 
\end{ex}

% nguồn: MaHoaNH_TracNghiem_Lop10
\begin{ex}%[0H5H4-1]%[ma:4759eccc]
	Cho hình vuông $ABCD$ tâm $O$. Số đo góc $\left(\overrightarrow{CO},  \overrightarrow{BA}\right)$ bằng
	\choice 
		{\True $45^{\circ}$} 
		{$145^{\circ}$} 
		 {$135^{\circ}$}
		 {$30^{\circ}$}
	\loigiai{ 
	\immini{Hình  vuông $ABCD$ có tâm $O$ nên
		 $$
		 \left(\overrightarrow{CO}, \overrightarrow{BA}\right)=\left(\overrightarrow{CO}, \overrightarrow{CD}\right)=\widehat{DCO}=45^{\circ}.
		$$}{\begin{tikzpicture}[scale=1,font=\footnotesize,line join=round,line cap=round,>=stealth]
		\path  (0,0) coordinate (B)
			(2,0) coordinate  (C)
			 ($(B)!1!90:(C)$) coordinate (A)
			($(A)+(C)-(B)$)  coordinate (D)
			($(A)!0.5!(C)$) coordinate (O); 
		 \draw (A)--(B)--(C)--(D)--cycle (A)--(C)
			 (B)--(D);
		\foreach \p/\g in {A/120,B/-150,C/-30,D/60,O/-90} \fill[black] (\p) circle (1pt)+(\g:0.3) node{$\p$}; 
	\end{tikzpicture}
	 }}
\end{ex}

% nguồn: MaHoaNH_TracNghiem_Lop10
\begin{ex}%[0H5H4-1]%[ma:8d820b0a]
	Cho tam giác $ABC$ đều cạnh bằng $a$,  trọng tâm $G$. Tích vô hướng của hai vectơ $\overrightarrow{BC}\cdot \overrightarrow{CG}$ bằng
	 \choice
		 {$\dfrac{a^2}{\sqrt{2}}$}
		 {$-\dfrac{a^2}{\sqrt{2}}$}
		{$\dfrac{a^2}{2}$} 
		{\True  $-\dfrac{a^2}{2}$}
	 \loigiai{
	\immini{Ta có \begin{align*}
		\left(\overrightarrow{BC},\overrightarrow{CG}\right)&=\left(-\overrightarrow{CB},\overrightarrow{CG}\right)=180^\circ-\left(\overrightarrow{CB},\overrightarrow{CG}\right)\\ 
		&=180^\circ-30^\circ=150^\circ 
	\end{align*}
	và $\left|\overrightarrow{CG}\right|=CG=\dfrac{2}{3}\cdot\dfrac{a\sqrt{3}}{2}=\dfrac{a\sqrt{3}}{3}$.}{\begin{tikzpicture}[scale=1,font=\footnotesize,line join=round,line cap=round,>=stealth]
		\path (0,0) coordinate (B)++(3,0) coordinate  (C)
			($(B)!1!60:(C)$) coordinate (A) 
			 (barycentric cs:A=1,B=1,C=1) coordinate (G)
			(barycentric  cs:B=1,A=1) coordinate (M);
		 \draw (A)--(B)--(C)--cycle (C)--(M);
		\foreach  \p/\g in {A/90,B/-150,C/-30,G/45,M/135} \fill[black] (\p) circle (1pt)+(\g:0.3) node{$\p$};
		\draw pic[draw,angle  radius=0.2cm]{right angle=B--M--C};
	\end{tikzpicture}
	} 
	\noindent\\ 
Vậy  $\overrightarrow{BC}\cdot\overrightarrow{CG}=\left|\overrightarrow{BC}\right|\cdot\left|\overrightarrow{CG}\right|\cdot\cos\left(\overrightarrow{BC},\overrightarrow{CG}\right)=a\cdot\dfrac{a\sqrt{3}}{3}\cdot\left(-\dfrac{\sqrt{3}}{2}\right)=-\dfrac{a^2}{2}$. }
\end{ex}

% nguồn: MaHoaNH_TracNghiem_Lop10
\begin{ex}%[0H5H4-1]%[ma:b86478d4]
	Cho tam giác $ABC$  vuông tại $A$ có $AB=a$, $AC=a\sqrt{3}$ và $AM$ là trung tuyến. Tính tích vô hướng $\overrightarrow{BA}\cdot \overrightarrow{AM}$.
	 \choice
		{$-a^2$} 
		 {$a^2$}
		{\True $-\dfrac{a^2}{2}$} 
		{$\dfrac{a^2}{2}$} 
	\loigiai{ 
	\immini{Ta có  $BC=\sqrt{AB^2+AC^2}=2a$. Suy ra $AM=BM=CM=a$.\\
		 Tam giác $ABM$ có $MA=MB=AB=a$ nên tam giác $ABM$ đều hay $\widehat{BAM}=60^\circ$.\\
		Khi đó \begin{align*}
		 \overrightarrow{BA}\cdot\overrightarrow{AM}&=-\overrightarrow{AB}\cdot\overrightarrow{AM}\\
		&=-AB\cdot  AM\cdot \cos \widehat{BAM}\\
		&=-a\cdot a\cdot  \dfrac{1}{2}\\
		&=-\dfrac{a^2}{2}. 
	\end{align*}}{\begin{tikzpicture}[scale=1,font=\footnotesize,line join=round,line cap=round,>=stealth]
		\path (0,0) coordinate  (B)++(3,0) coordinate (C)
			 ($(B)!0.5!(C)$) coordinate (M)
			 ($(M)!1!-60:(B)$) coordinate (A);
		\draw (A)--(B)--(C)--cycle (A)--(M); 
		\foreach \p/\g in {A/90,C/-30,B/210,M/-90} \fill[black] (\p)  circle (1pt)+(\g:0.3) node{$\p$};
		\draw pic[draw,angle radius=0.2cm]{right angle=B--A--C}; 
	\end{tikzpicture}
}} 
\end{ex}

% nguồn: MaHoaNH_TracNghiem_Lop10
\begin{ex}%[0H5H4-1]%[ma:1c82ad0c]
	Cho hình chữ nhật $ABCD$ có $AB=a, AC=2 a$. Tính  góc giữa hai vectơ $\overrightarrow{CA}$ và $\overrightarrow{DC}$.
	 \choice
		{$60^{\circ}$} 
		{$45^{\circ}$} 
		 {$150^{\circ}$}
		{\True $120^{\circ}$} 
	\loigiai{ 
	\immini{Ta có \begin{align*}
		 \left(\overrightarrow{CA},\overrightarrow{DC}\right)&=\left(\overrightarrow{CA},-\overrightarrow{CD}\right)\\
		&=180^\circ-\left(\overrightarrow{CA},\overrightarrow{CD}\right)=180^\circ-\widehat{DCA}. 
	\end{align*}
	}{\begin{tikzpicture}[scale=1,font=\footnotesize,line join=round,line cap=round,>=stealth]
		\path (0,0) coordinate (b)++(3,0) coordinate  (C)++(0,1.5) coordinate (D)
			($(b)+(D)-(C)$) coordinate (a) 
			($(C)!1!-30:(b)$) coordinate  (c)
			(intersection of D--a and C--c) coordinate  (A)
			 ($(b)!(A)!(C)$) coordinate (B);
		 \draw (A)--(B)--(C)--(D)--cycle (A)--(C);
		\foreach \p/\g in  {A/120,B/-150,C/-30,D/60} \fill[black] (\p) circle (1pt)+(\g:0.3) node{$\p$};
		\foreach \x/\y/\z in {D/A/B,A/B/C,B/C/D,C/D/A} \draw pic[draw,angle radius=0.2cm]{right  angle=\x--\y--\z};
	\end{tikzpicture}
	}\noindent\\ 
	Tam giác $ACD$  vuông tại $D$ nên $\cos\widehat{DCA}=\dfrac{CD}{CA}=\dfrac{a}{2a}=\dfrac{1}{2}\Rightarrow\widehat{DCA}=60^\circ$.\\
	Vậy  $\left(\overrightarrow{CA},\overrightarrow{DC}\right)=180^\circ-60^\circ=120^\circ$.
 }
\end{ex}

% nguồn: MaHoaNH_TracNghiem_Lop10
\begin{ex}%[0H5H4-1]%[ma:2462a636]
	Cho tam giác $ABC$ đều. Giá trị $\sin \left( \overrightarrow{BC},\overrightarrow{AC}  \right)$ là
	 \choice
		 {$\dfrac{1}{2}$}
		 {$-\dfrac{1}{2}$}
		{$-\dfrac{\sqrt 3}  2$}
		{\True $\dfrac{\sqrt 3}  2$}
	\loigiai{ 
	\immini{ 
		Ta có  $\left( \overrightarrow{BC},\overrightarrow{AC} \right)$ $=\left( \overrightarrow{CB'},\overrightarrow{CA'} \right)$ $=\widehat{A'CB}$ $=60^\circ$.\\
		Vậy $ \sin \left( \overrightarrow{BC},\overrightarrow{AC} \right)$  $=\sin 60^\circ $ $=\dfrac{\sqrt 3} 2$.}{
		\begin{tikzpicture}[line join = round,line cap = round,scale=1]
				 \def \a{2};
				\coordinate (A) at  (0,0);
				\coordinate (C)  at ($(A)+(\a,0)$);
				\coordinate[shift=(60:\a  cm)] (B) at (A);
				\coordinate (B') at ($(C)+(C)-(B)$); 
				\coordinate (A') at ($(C)+(C)-(A)$); 
				 \draw pic[draw,angle eccentricity=1.4,pic text=$ $,angle radius=0.35cm]{angle=B'--C--A'};
				 \draw (A)--(B)--(C)--cycle;
				 \draw[->,arrows={-latex}] (C)--(A');
				 \draw[->,arrows={-latex}] (C)--(B');
				\foreach  \x/\y in {A/180,B/90,C/30,B'/-90,A'/0}
				 \draw[fill=black] (\x) circle (0.7pt) + (\y:0.5cm) node{$\x$};
			\end{tikzpicture}
		 }
	} 
\end{ex}

% nguồn: MaHoaNH_TracNghiem_Lop10
\begin{ex}%[0H5H4-1]%[ma:db0a9bd8]
	Cho tam giác $ABC$ vuông cân tại $A$, $AB=1$. Khẳng định nào  sau đây \textbf{sai}?
	 \choice
		{$\overrightarrow{AB}  \cdot \overrightarrow{BC}=-1$}
		{$\overrightarrow{CA}  \cdot \overrightarrow{CB}=1$}
		{$\overrightarrow{AB}  \cdot \overrightarrow{AC}=0$}
		{\True  $\overrightarrow{AB} \cdot \overrightarrow{CB}=-1$}
	 \loigiai{
	 \immini{
		Dựng  $\overrightarrow{AD}=\overrightarrow{CB}$.\\
		Ta  có
		$\overrightarrow{AB} \cdot \overrightarrow{CB}=\overrightarrow{AB} \cdot  \overrightarrow{AD}=AB\cdot AD\cdot \cos{45^\circ}=\sqrt 2\cdot \dfrac{\sqrt 2} 2=1$.}{
		\begin{tikzpicture}[line join = round,line cap = round,scale = 0.6]
				 \def\a{3}
				 \coordinate[label = left:$A$] (A) at (0,0);
				\coordinate[label = right:$C$]  (C) at (\a,0);
				\coordinate[label  = above:$B$] (B) at (0,\a);
				\coordinate[label  = above:$D$] (D) at (-\a,\a);
				\foreach \x/\y in  {A/C,A/B}{
					\path  (\x)--(\y) node[midway,sloped]{\tikz{\draw (-90:2pt)--(90:2pt);}};
					\draw  (\x)--(\y) node[midway,sloped] {};
				} 
				\draw[fill  = gray!50] ($(A)!8pt!(B)$)--($(A)!8pt!(B)+(A)!8pt!(C)-(A)$)--($(A)!8pt!(C)$)--(A)--cycle;
				 \draw (A)--(C)--(B)--(D)--(A)--cycle;
				\foreach \x in {A,B,C} \fill[black] (\x) circle  (1.5pt);
			\end{tikzpicture}
		} 
	 }
\end{ex}

% nguồn: MaHoaNH_TracNghiem_Lop10
\begin{ex}%[0H5H4-1]%[ma:6dd909b0]
	Cho tam giác  $ABC$ vuông tại $A$ có $\widehat{B}=30^\circ,AC=2.$ Gọi $M$ là trung điểm của $BC.$ Tính giá trị của biểu thức $P=\overrightarrow{AM} \cdot \overrightarrow{BM}.$
	\choice 
		{\True $P=-2$} 
		{$P=2\sqrt 3$} 
		{$P=2$} 
		{$P=-2\sqrt 3$} 
	 \loigiai{
	 \immini{
		Tam giác $ABC$ vuông tại $A$  nên
		\[AB=\dfrac{AC}{\tan B}=\dfrac  2{\tan 30^\circ}=2\sqrt 3.\]
		Suy ra $BC=\sqrt{AB^2+AC^2}=\sqrt{\left( 2\sqrt 3 \right)^2+2^2}=4.$ 
		Khi đó \begin{eqnarray*}
		P&=&\overrightarrow{AM}  \cdot \overrightarrow{BM}=\dfrac{1}{2}\left( \overrightarrow{AB}+\overrightarrow{AC} \right)\cdot \dfrac{1}{2}\overrightarrow{BC}\\
		&=&\dfrac{1}{4}\left( \overrightarrow{AB} \cdot \overrightarrow{BC}+\overrightarrow{AC} \cdot  \overrightarrow{BC} \right)\\
		&=&\dfrac{1}{4}\left( -AB\cdot BC\cdot \cos  B+AC\cdot BC\cdot \cos C \right)\\
		&=&\dfrac{1}{4}\left( -2\sqrt 3\cdot 4\cdot \cos 30^\circ+2\cdot 4\cdot \cos 60^\circ  \right)=-2.
	\end{eqnarray*}
	 }{
	\begin{tikzpicture}[scale=1,font=\footnotesize,line join=round,line cap=round,>=stealth]
			 \path
				 (-2,0) coordinate (B)
				(2,0) coordinate  (C)
				(60:2) coordinate (A) 
				 ($(B)!0.5!(C)$) coordinate (M);;
			\draw (A)--(B)--(C)--(A)--(M); 
			\draw pic[draw,angle eccentricity=1.4,pic text=$30^\circ  $,angle radius=0.75cm]{angle=C--B--A};
			\draw pic[draw,angle radius = 6pt]{right angle =  B--A--C};
			\foreach \p/\q in {A/90,B/-90,C/-90,M/-90} \fill[black] (\p) circle (1pt) 
				($(\p)+(\q:2.5mm)$) node{$\p$}; 
		\end{tikzpicture}
 }}
\end{ex}

% nguồn: MaHoaNH_TracNghiem_Lop10
\begin{ex}%[0H5H4-1]%[ma:cdb48859]
	Cho tam giác $ABC$ đều có trọng tâm $G$ và  $H$ là trung điểm của $BC$. Xác định $\cos \left( \overrightarrow{GB},\overrightarrow{GH} \right)$.
	 \choice
		{\True $\dfrac{1}{2}$} 
		 {$-\dfrac{1}{2}$}
		{$-\dfrac{\sqrt  3} 2$}
		{$\dfrac{\sqrt  3} 2$}
	\loigiai{ 
	\immini{ 
		Do tam giác $ABC$ đều nên  $AH\perp BC$ và $\widehat{GBH}=30^\circ $. \\
		Suy  ra $\left( \overrightarrow{GB},\overrightarrow{GH} \right)=\widehat{BGH}=60^\circ $.\\
		Vậy $\cos  \left( \overrightarrow{GB},\overrightarrow{GH} \right)=\cos 60^\circ =\dfrac{1}{2}$.}{
		\begin{tikzpicture}[line join = round,line cap = round,scale=1]
				\def \a{3}; 
				 \coordinate (B) at (0,0);
				\coordinate (C)  at ($(B)+(\a,0)$);
				\coordinate[shift=(60:\a cm)]  (A) at (B);
				\coordinate (H) at  ($(C)!0.5!(B)$);
				\coordinate (G) at ($(A)!2/3!(H)$); 
				\draw pic[draw,angle eccentricity=1.4,pic text=$  $,angle radius=0.35cm]{angle=B--G--H};
				\draw (A)--(B)--(C)--cycle (B)--(G) 
					 (A)--(H);
				 \draw[->,arrows={-latex}] (G)--(B);
				\draw[->,arrows={-latex}] (A)--(H); 
				\foreach \x/\y  in {A/C,A/B,B/C}{
					 \path (\x)--(\y) node[midway,sloped]{\tikz{\draw (-90:2pt)--(90:2pt);}};
					 \draw (\x)--(\y) node[midway,sloped] {};
				} 
				\foreach  \x/\y in {A/90,B/180,C/0,H/-90,G/30}
				\draw[fill=black] (\x) circle (0.7pt) +  (\y:0.25cm) node{$\x$};
			\end{tikzpicture}
		} 
	 }
\end{ex}

% nguồn: MaHoaNH_TracNghiem_Lop10
\begin{ex}%[0H5H4-1]%[ma:7db7a073]
	Cho tam giác $ABC$ vuông tại $A$ có $\widehat{B}=30^\circ$, $AC=2$.  Gọi $M$ là trung điểm của $BC$. Giá trị của biểu thức $P=\vv{AM}\cdot \vv{BM}$ bằng
	 \choice
		 {\True $-2$}
		 {$2\sqrt{3}$}
		{ $2$} 
		{$-2\sqrt{3}$} 
	\loigiai{ 
		\begin{center}
			\begin{tikzpicture}[scale=1,line join=round,font=\footnotesize,line cap=round,thick]
					\coordinate (B) at  (-3,0);
					\coordinate (C) at (2,0); 
					\coordinate (Temp)  at ($(B) + (30:1)$);
					\coordinate (A) at  ($(B)!(C)!(Temp)$);
					\path ($(B)!0.5!(C)$) coordinate  (M);
					 \path ($(A)!0.5!(C)$) node [above right]{$2$};
					\draw (A)--(B)--(C)--cycle (A)--(M); 
					\pic[draw,thin,angle  radius=2mm] {right angle = B--A--C};
					 \foreach \i/\g in {A/90,B/-90,C/-90,M/-90}{\draw[fill=black](\i) circle (1pt)
						($(\i)+(\g:3mm)$)  node[scale=1]{$\i$};}
					 \pic["$30^\circ$",draw,angle radius=5mm,angle eccentricity=1.6]{angle=C--B--A};
				\end{tikzpicture}
		\end{center}
		Xét tam giác $ABC$  vuông tại $A$ có $\widehat{B}=30^\circ\Rightarrow \widehat{C}=60^\circ$.\\
		Mặt khác $M$ là trung  điểm của $BC$ nên $AM=MC=\dfrac{BC}{2}$.\\
		Do đó tam giác $MAC$ đều $\Rightarrow  MA=MC=AC=2$.\\
		 Khi đó
		\begin{align*}
			P=\vv{AM}\cdot \vv{BM}&=\vv{AM}\cdot  \vv{MC}\\
			 &=-\vv{MA}\cdot \vv{MC}\\
			&=-\left|\vv{MA} \right|\cdot \left|\vv{MC}\right|\cdot  \cos \left(\vv{MA},\vv{MC}\right)\\
			&=-2\cdot 2\cdot \cos 60^\circ=-2. 
		\end{align*}
	 }
\end{ex}

% nguồn: MaHoaNH_TracNghiem_Lop10
\begin{ex}%[0H5V4-1]%[ma:6d99a858]
	Cho tam giác $ABC$ vuông tại $A$, biết $\widehat{B}=50^{\circ}$. Kẻ đường cao  $AH$ $(H \in BC)$, đường phân giác trong của góc $C$ là $CK$ $(K \in AB)$. Xác định góc giữa hai vectơ $\overrightarrow{AH}$ và $\overrightarrow{CK}$.
	\choice 
		 {\True $110^{\circ}$}
		 {$120^{\circ}$}
		 {$100^{\circ}$}
		{$90^{\circ}$} 
	\loigiai{ 
	\immini{Gọi $D$ là giao  điểm của $AH$ và $CK$.\\
		Ta có $\overrightarrow{AH}$ cùng hướng với  $\overrightarrow{DH}$ và $\overrightarrow{CK}$ cùng hướng với $\overrightarrow{DK}$ nên \[\left(\overrightarrow{AH},\overrightarrow{CK}\right)=\left(\overrightarrow{DH},\overrightarrow{DK}\right)=\widehat{HDK}.\]
		Sử dụng tính chất góc ngoài tam giác ta có $$\widehat{BKC}=\widehat{KAC}+\widehat{ACK}=90^\circ+20^\circ=110^\circ.$$}{\begin{tikzpicture}[scale=1,font=\footnotesize,line join=round,line cap=round,>=stealth]
		 \def\r{2};
		\path (0,0) coordinate (O)++(0:{\r}) coordinate (C) 
			($(C)!2!(O)$) coordinate  (B)
			($(B)!1!50:(C)$) coordinate  (b)
			($(C)!1!-40:(B)$)  coordinate (c)
			(intersection of B--b and C--c) coordinate (A) 
			($(C)!1!-20:(B)$) coordinate  (k)
			(intersection of B--A and C--k)  coordinate (K)
			($(B)!(A)!(C)$) coordinate (H) 
			(intersection of A--H and C--K) coordinate (D); 
		\draw (A)--(B)--(C)--cycle (C)--(K) 
			(A)--(H); 
		 \foreach \p/\g in {A/90,B/-150,C/-30,K/150,H/-90,D/45} \fill[black] (\p) circle (1pt)+(\g:0.3) node{$\p$};
		\foreach \x/\y/\z in {C/A/B,A/H/C} \draw pic[draw,angle radius=0.2cm]{right  angle=\x--\y--\z};
	\end{tikzpicture}
		 }\noindent\\
		Vậy $\widehat{HDK}=360^\circ-\widehat{DHB}-\widehat{HBK}-\widehat{BKD}=360^\circ-90^\circ-50^\circ-110^\circ=110^\circ$. 
	} 
\end{ex}


%%%%%%%%%%%%%%%% ĐÚNG/SAI %%%%%%%%%%%%%%%%

% nguồn: MaHoaNH_DungSai_Lop10
\begin{ex}%[0H5N4-1]%[ma:51379eec]
	Biết rằng $\left|\overrightarrow{a}\right|=3$, $\left|\overrightarrow{b}\right|=4$ và $\overrightarrow{a}\cdot\overrightarrow{b}=12$. Lúc đó, $\alpha=  \left(\overrightarrow{a},\overrightarrow{b}\right)$ thỏa mãn:
	 \choiceTF
		{\True $\alpha  = 0^\circ$}
		 {\True $\cos \alpha = 1$}
		{$\alpha = 90^\circ$} 
		 {$\alpha = 45^\circ$}
	 \loigiai{
		Ta có $\cos \left(\overrightarrow{a},\overrightarrow{b}\right)=\dfrac{\overrightarrow{a}\cdot\overrightarrow{b}}{\left|\overrightarrow{a}\right|\cdot\left|\overrightarrow{b}\right|}=\dfrac{12}{3\cdot 4}=1$. \\ 
		 Do đó, $\left(\overrightarrow{a},\overrightarrow{b}\right)=0^\circ$.
	 }
\end{ex}

% nguồn: MaHoaNH_DungSai_Lop10
\begin{ex}%[0H5N4-1]%[ma:229ef3d4]
	Cho đoạn  thẳng $AB$. Điểm $M$ thuộc đường tròn đường kính $AB$.
	Xét  tính đúng sai của các mệnh đề sau
	 \choiceTF
		{\True  $\overrightarrow{MA}\cdot\overrightarrow{MB}=0$}
		 {\True $\overrightarrow{MB}\cdot\overrightarrow{MA}=0$}
		 {$\overrightarrow{MA}\cdot\overrightarrow{MB}=1$}
		 {$\overrightarrow{MB}\cdot\overrightarrow{MA}=1$}
	 \loigiai{
		$M$ thuộc đường tròn đường kính $AB$ nên $MA \perp MB$.\\ 
		Nên $\overrightarrow{MA}\cdot \overrightarrow{MB}=0$  và $\overrightarrow{MB}\cdot \overrightarrow{MA}=0$
	} 
\end{ex}

% nguồn: MaHoaNH_DungSai_Lop10
\begin{ex}%[0H5N4-1]%[ma:b075f0cb]
	Cho tam giác đều $A B  C$, có độ dài các cạnh bằng $1$. Xác định tính đúng sai của các mệnh đề sau
	 \choiceTF
		 {\True $\left(\overrightarrow{A B}; \overrightarrow{A C} \right)=\widehat{CAB}$}
		{\True  $\overrightarrow{A B} \cdot \overrightarrow{A C}=\dfrac{1}{2}$}
		{$\left(\overrightarrow{A B}; \overrightarrow{A C}  \right)=-60^{\circ}$}
		{$\overrightarrow{A B}  \cdot \overrightarrow{A C}=-\dfrac{1}{2}$}
	 \loigiai{
		Do tam giác $ABC$ đều, nên $\widehat{C A B}=\widehat{ABC}=\widehat{BCA}=60^{\circ}$.  \\
		Suy ra $\left(\overrightarrow{A B}; \overrightarrow{A  C} \right)=\widehat{CAB}=60^{\circ}$. \\
		$\overrightarrow{A B} \cdot \overrightarrow{A C}=|\overrightarrow{A B}| \cdot|\overrightarrow{A C}| \cdot \cos \left(\overrightarrow{A B}; \overrightarrow{A C} \right)=1 \cdot 1 \cdot \cos 60^{\circ}=\dfrac{1}{2}$. 
	} 
\end{ex}

% nguồn: MaHoaNH_DungSai_Lop10
\begin{ex}%[0H5N4-1]%[ma:7fcbedde]
	Cho bốn điểm $A$, $B$, $C$, $D$. Xác định  tính đúng sai của các mệnh đề sau
	\choiceTF 
		{\True $\overrightarrow{AB}\cdot(\overrightarrow{AD}-\overrightarrow{AC})+\overrightarrow{AC}\cdot(\overrightarrow{AB}-\overrightarrow{AD})+\overrightarrow{AD}\cdot(\overrightarrow{AC}-\overrightarrow{AB})=0$} 
		{\True 
		$\overrightarrow{AB}\cdot\overrightarrow{CD}+\overrightarrow{AC}\cdot\overrightarrow{DB}+\overrightarrow{AD}\cdot\overrightarrow{BC}=0$} 
		{ $\overrightarrow{AB}\cdot(\overrightarrow{AD}-\overrightarrow{AC})+\overrightarrow{AC}\cdot(\overrightarrow{AB}-\overrightarrow{AD})+\overrightarrow{AD}\cdot(\overrightarrow{AC}-\overrightarrow{AB})=AB$} 
		{  $\overrightarrow{AB}\cdot\overrightarrow{CD}+\overrightarrow{AC}\cdot\overrightarrow{DB}+\overrightarrow{AD}\cdot\overrightarrow{BC}=AB^2$}
	 \loigiai{
		 Ta có
		\begin{eqnarray*}
			& & \overrightarrow{AB}\cdot\overrightarrow{CD}+\overrightarrow{AC}\cdot\overrightarrow{DB}+\overrightarrow{AD}\cdot\overrightarrow{BC} \\ 
			& =  & \overrightarrow{AB}\cdot(\overrightarrow{AD}-\overrightarrow{AC})+\overrightarrow{AC}\cdot(\overrightarrow{AB}-\overrightarrow{AD})+\overrightarrow{AD}\cdot(\overrightarrow{AC}-\overrightarrow{AB})=0.
		\end{eqnarray*}
	} 
\end{ex}

% nguồn: MaHoaNH_DungSai_Lop10
\begin{ex}%[0H5N4-1]%[ma:e9144f6a]
	Cho tam  giác $ABC$. Giá trị của biểu thức $\overrightarrow{BA}\cdot\overrightarrow{CA}$ bằng
	\choiceTF 
		 {\True $AB\cdot AC\cdot\cos\widehat{BAC}$}
		{$-AB\cdot  AC\cdot\cos\widehat{BAC}$}
		 {$AB\cdot AC\cdot\cos\widehat{ABC}$}
		 {$AB\cdot AC\cdot\cos\widehat{ACB}$}
	\loigiai{ 
		Ta có 
		$\begin{aligned}
		\overrightarrow{BA}\cdot\overrightarrow{CA} 
		&=\left(-\overrightarrow{AB}\right)\cdot\left(-\overrightarrow{AC}\right) 
		 =\overrightarrow{AB}\cdot\overrightarrow{AC}\\
		 &=AB\cdot AC\cdot\cos(\widehat{BAC}).
	\end{aligned}$
 }
\end{ex}

% nguồn: MaHoaNH_DungSai_Lop10
\begin{ex}%[0H5N4-1]%[ma:32aab9cc]
	Cho đoạn thẳng $AB$. Biết $\overrightarrow{MA}\cdot\overrightarrow{MB}=0$. Lúc  đó, điểm $M$ có thể:
	 \choiceTF
		{\True Trùng  với điểm $A$}
		{\True  Trùng với điểm $B$}
		{Trùng với trung điểm $I$  của $AB$}
		 {\True Thuộc đường tròn đường kính $AB$}
	 \loigiai{
		\begin{itemchoice}
			\itemch $\overrightarrow{AA}\cdot\overrightarrow{AB}=0$ 
			 \itemch $\overrightarrow{BA}\cdot\overrightarrow{BB}=0$
			\itemch  $\overrightarrow{IA}\cdot\overrightarrow{IB}=-IA^2$
			\itemch  $\overrightarrow{MA}\cdot\overrightarrow{MB}=0$ nên $MA \perp MB$.\\
			 Vậy $M$ thuộc đường tròn đường kính $AB$.
		\end{itemchoice}
	 }
\end{ex}

% nguồn: MaHoaNH_DungSai_Lop10
\begin{ex}%[0H5H4-1]%[ma:f5f7cbf8]
	Cho tam giác $ABC$ vuông tại $A$  có $AB=a, BC=2 a$. Khi đó
	\choiceTF 
		 {\True $\widehat{ACB}=60^{\circ}$}
		{\True  $\overrightarrow{BA}\cdot \overrightarrow{BC}=a^2$}
		{$\overrightarrow{BC}\cdot  \overrightarrow{CA}=3 a^2$}
		{\True $\overrightarrow{AB}\cdot \overrightarrow{BC}+\overrightarrow{BC}\cdot \overrightarrow{CA}+\overrightarrow{CA}\cdot  \overrightarrow{AB}=-4 a^2$}
	\loigiai{\begin{center}
		\begin{tikzpicture}[scale=1,font=\footnotesize,line join=round,line cap=round,>=stealth]
				\path (0,0) coordinate (A)++(3,0) coordinate  (C)
					 (A)++(0,2) coordinate (B);
				 \draw (A)--(B)--(C)--cycle;
				\foreach \p/\g in {A/210,B/90,C/-30} \fill[black] (\p) circle  (1pt)+(\g:0.3) node{$\p$};
				\draw pic[draw,angle radius=0.2cm]{right angle=C--A--B}; 
			\end{tikzpicture}
	\end{center}
	\begin{itemchoice}
		\itemch Đúng.  Ta có $\sin\widehat{ACB}=\dfrac{AB}{BC}=\dfrac{1}{2}\Rightarrow\widehat{ACB}=60^\circ$.
		\itemch Đúng. Ta  có $AC=\sqrt{BC^2-AB^2}=a\sqrt{3}$. Khi đó, \[\overrightarrow{BA}\cdot\overrightarrow{BC}=\dfrac{1}{2}\left(BA^2+BC^2-AC^2\right)=\dfrac{1}{2}\left(a^2+4a^2-3a^2\right)=a^2.\]
		\itemch Sai. Ta có \begin{align*}
		\overrightarrow{BC}\cdot \overrightarrow{CA}=-\overrightarrow{CB}\cdot\overrightarrow{CA}=-CB\cdot  CA\cdot\cos\widehat{ACB}=-2a\cdot a\cdot \dfrac{1}{2}=-a^2.
	\end{align*}
	\itemch Đúng. Ta có \begin{eqnarray*}
	&&\overrightarrow{AB}\cdot  \overrightarrow{BC}+\overrightarrow{BC}\cdot \overrightarrow{CA}+\overrightarrow{CA}\cdot \overrightarrow{AB}\\
	 &=&-\overrightarrow{BA}\cdot \overrightarrow{BC}-\overrightarrow{CB}\cdot \overrightarrow{CA}-\overrightarrow{AC}\cdot \overrightarrow{AB}\\
	 &=&-\dfrac{1}{2}\left(BA^2+BC^2-AC^2\right)-\dfrac{1}{2}\left(CB^2+CA^2-AB^2\right)-\dfrac{1}{2}\left(AC^2+AB^2-BC^2\right)\\
	&=&-\dfrac{1}{2}\left(AB^2+BC^2+CA^2\right)\\ 
	&=&-\dfrac{1}{2}\left(a^2+4a^2+3a^2\right)\\ 
	&=&-4a^2. 
	\end{eqnarray*}
	\end{itemchoice}
} 
\end{ex}

% nguồn: MaHoaNH_DungSai_Lop10
\begin{ex}%[0H5H4-1]%[ma:d38bc702]
	 \immini[thm]{
	Cho hình chữ nhật $ABCD, AB=4 a, AD=3 a$. Gọi $M$ là trung điểm  của $AB$, $G$ là trọng tâm tam giác $ACM$ (xem hình bên). Khi đó
	 \choiceTF
		 {$\overrightarrow{CM}=\dfrac{1}{2}\overrightarrow{BA}-3\overrightarrow{BC}$}
		{$\overrightarrow{BG}=\dfrac{3}{2}\overrightarrow{BA}+\dfrac{1}{3}\overrightarrow{BC}$} 
		{\True $\overrightarrow{BC}\cdot \overrightarrow{BA}=0$} 
		{$\overrightarrow{BG}\cdot  \overrightarrow{CM}=-a^2$}}{
		\begin{tikzpicture}[line join = round,line cap = round,>=stealth,font=\footnotesize,scale=1]
				 \coordinate (D) at ($(0,0)$);
				\coordinate (C) at ($(D)+(5,0)$); 
				\coordinate (A) at ($(D)+(0,3)$); 
				\coordinate (B)  at ($(A)+(C)-(D)$);
				 \coordinate (M) at ($(A)!0.5!(B)$);
				 \coordinate (N) at ($(A)!0.5!(M)$);
				\coordinate  (G) at ($(C)!2/3!(N)$);
				\draw  (A)--(B)--(C)--(D)--(A)--(C)--(M)
					(G)--(B)--(D); 
				\foreach \x/\y  in {A/90,B/90,D/-90,C/-90,M/90,G/90}
				\draw[fill=black]  (\x) circle (1pt) + (\y:0.25cm) node{$\x$};
			\end{tikzpicture}
		} 
	 \loigiai{
		\begin{itemchoice}
			 \itemch Sai. $\overrightarrow{CM}=\overrightarrow{BM}-\overrightarrow{B C}=\dfrac{1}{2} \overrightarrow{BA}-\overrightarrow{BC}$.
			\itemch Sai. Vì $G$ là  trọng tâm của tam giác ${ACM}$ nên\\
			 $3 \overrightarrow{BG}=\overrightarrow{BA}+\overrightarrow{BM}+\overrightarrow{BC}=\overrightarrow{BA}+\dfrac{1}{2} \overrightarrow{BA}+\overrightarrow{BC}=\dfrac{3}{2} \overrightarrow{BA}+\overrightarrow{BC} \Rightarrow \overrightarrow{BG}=\dfrac{1}{2} \overrightarrow{BA}+\dfrac{1}{3} \overrightarrow{BC}$.
			\itemch Đúng. Vì $ABCD$ là hình chữ nhật  nên $BC \perp BA \Rightarrow \overrightarrow{BC} \cdot \overrightarrow{BA}=0$.
			\itemch Sai. Vì \begin{eqnarray*}
			\overrightarrow{BG} \cdot  \overrightarrow{CM}
			 &=&\left(\dfrac{1}{2} \overrightarrow{BA}+\dfrac{1}{3} \overrightarrow{BC}\right) \cdot\left(\dfrac{1}{2} \overrightarrow{BA}-\overrightarrow{BC}\right)\\
			&=&\dfrac{1}{4}  \overrightarrow{B A}^2-\dfrac{1}{3} \overrightarrow{BA} \cdot \overrightarrow{BC}-\dfrac{1}{3} \overrightarrow{BC}^2\\
			&=&\dfrac{1}{4}{{(4a)}^2}-\dfrac{1}{3}\cdot  4a\cdot 3a-\dfrac{1}{3}{{(3a)}^2}\\
			&=&-3a^2. 
		\end{eqnarray*}
	\end{itemchoice}
 }
\end{ex}

% nguồn: MaHoaNH_DungSai_Lop10
\begin{ex}%[0H5H4-1]%[ma:4dc5135a]
	\immini[thm]{Cho hình vuông $A B C D$ cạnh $ a$.  Lấy $E$ là trung điểm của ${B C}$, điểm $F$ thoả mãn $\overrightarrow{B F}=\dfrac{3}{4} \overrightarrow{B D}$. Khi đó
	 \choiceTF
		{\True  $\overrightarrow{AE}=\overrightarrow{AB}+\dfrac{1}{2}\overrightarrow{AD}$}
		{$\overrightarrow{AF}=\dfrac{1}{4}\overrightarrow{AB}+\dfrac{5}{4}\overrightarrow{AD}$} 
		{\True $\overrightarrow{EF}=\dfrac{-3} 4\overrightarrow{AB}+\dfrac{1}{4}\overrightarrow{AD}$} 
		{\True Tam giác $A E F$  vuông cân}}{
		\begin{tikzpicture}[line join = round,line cap = round,>=stealth,font=\footnotesize,scale=1]
				\coordinate (A)  at ($(0,0)$);
				\coordinate (B) at  ($(A)+(3,0)$);
				\coordinate (D) at  ($(A)+(0,3)$);
				\coordinate (C) at ($(A)+(B)+(D)$); 
				\coordinate (E)  at ($(B)!0.5!(C)$);
				\coordinate (F) at ($(B)!3/4!(D)$); 
				\draw (A)--(B)--(C)--(D)--(A)--(E)--(F)--(A) 
					 (B)--(D);
				 \foreach \x/\y in {A/-90,B/-90,D/90,C/90,E/0,F/90}
				\fill (\x) circle (1pt) + (\y:0.25cm) node{$\x$}; 
			\end{tikzpicture}
		} 
	 \loigiai{
		\begin{itemchoice}
			\itemch Đúng. Vì $\overrightarrow{A E}=\overrightarrow{A B}+\overrightarrow{B E}=\overrightarrow{A B}+\dfrac{1}{2} \overrightarrow{B C}=\overrightarrow{A B}+\dfrac{1}{2} \overrightarrow{A  D}$.
			\itemch Sai. Vì $\overrightarrow{AF}=\overrightarrow{AB}+\overrightarrow{BF}=\overrightarrow{AB}+\dfrac{3}{4}\overrightarrow{BD}=\overrightarrow{AB}+\dfrac{3}{4}\left(\overrightarrow{AD}-\overrightarrow{AB}\right)=\dfrac{1}{4}\overrightarrow{AB}+\dfrac{3}{4}\overrightarrow{AD}$. 
			\itemch Đúng. Vì $\overrightarrow{EF}=\overrightarrow{AF}-\overrightarrow{AE}=\left( \dfrac{1}{4}\overrightarrow{AB}+\dfrac{3}{4}\overrightarrow{AD} \right)-\left( \overrightarrow{AB}+\dfrac{1}{2}\overrightarrow{AD} \right)=\dfrac{-3}  4\overrightarrow{AB}+\dfrac{1}{4}\overrightarrow{AD}.$
			\itemch Đúng. Ta có 
			\begin{eqnarray*}
				\overrightarrow{A F} \cdot  \overrightarrow{E F}
				&=&\left(\dfrac{1}{4} \overrightarrow{A B}+\dfrac{3}{4} \overrightarrow{A D}\right) \cdot\left(\dfrac{-3} 4 \overrightarrow{A B}+\dfrac{1}{4}  \overrightarrow{A D}\right)\\
				 &=&\dfrac{-3}{16}{{\overrightarrow{AB}}^2}-\dfrac{1}{2}\overrightarrow{AB}\cdot \overrightarrow{AD}+\dfrac 3{16}{{\overrightarrow{AD}}^2}\\
				 &=&0.
			\end{eqnarray*}
			Vậy $ AF\perp EF$. \\ 
			Ta có \begin{eqnarray*}
			 \overrightarrow{A F}^2&=&\left(\dfrac{1}{4} \overrightarrow{A B}+\dfrac{3}{4} \overrightarrow{A D}\right)^2=\dfrac 1{16} \overrightarrow{A B}^2+\dfrac{3}{8} \overrightarrow{A B} \cdot \overrightarrow{A D}+\dfrac 9{16} \overrightarrow{A D}^2=\dfrac{5}{8} \overrightarrow{A B}^2.\\
			\overrightarrow{EF}^2&=&{{\left( \dfrac{-3}  4\overrightarrow{AB}+\dfrac{1}{4}\overrightarrow{AD} \right)}^2}=\dfrac 9{16}{{\overrightarrow{AB}}^2}-\dfrac{3}{8}\overrightarrow{AB}\cdot \overrightarrow{AD}+\dfrac 1{16}{{\overrightarrow{AD}}^2}=\dfrac{5}{8}{\overrightarrow{AB}}^2.
		\end{eqnarray*}
		Do đó, $A F^2=E F^2=\dfrac{5}{8} \overrightarrow{A B}^2 \Rightarrow A F=E F$. Vậy tam giác ${A E F}$ vuông cân tại ${F}$.\\ 
		Chú ý: Ta có thể chứng minh tam giác ${A E F}$ vuông bằng định  lí Pythagore.
	\end{itemchoice}
} 
\end{ex}

% nguồn: MaHoaNH_DungSai_Lop10
\begin{ex}%[0H5V4-1]%[ma:f2d916da]
	Cho hình  vuông $ABCD$ tâm $O$, có cạnh $a$. Biết $M$ là trung điểm của $AB, G$ là trọng tâm tam giác $ADM$. Khi đó
	\choiceTF 
		 {$\overrightarrow{AB}\cdot \overrightarrow{CA}=a^2$}
		{$\overrightarrow{AM}\cdot \overrightarrow{AC}=\dfrac{a^2}{3}$} 
		 {\True $\overrightarrow{AD}\cdot \overrightarrow{BD}+\overrightarrow{O M}\cdot \overrightarrow{AC}=\dfrac{a^2}{2}$}
		 {\True $\left(\overrightarrow{AB}+\overrightarrow{AD}\right)\left(\overrightarrow{BD}+\overrightarrow{BC}\right)=a^2$}
	 \loigiai{
		\begin{center}
			\begin{tikzpicture}[scale=1,font=\footnotesize,line join=round,line cap=round,>=stealth]
					\path  (0,0) coordinate (B)
						(3,0) coordinate  (C)
						($(B)!1!90:(C)$)  coordinate (A)
						 ($(A)+(C)-(B)$) coordinate (D)
						($(A)!0.5!(C)$) coordinate (O) 
						(barycentric  cs:A=1,B=1) coordinate (M)
						 (barycentric cs:A=1,D=1,M=1) coordinate (G);
					\draw  (A)--(B)--(C)--(D)--cycle (A)--(C)
						 (B)--(D)--(M);
					\foreach \p/\g in {A/120,B/-150,C/-30,D/60,M/180,G/0,O/-90} \fill[black]  (\p) circle (1pt)+(\g:0.3) node{$\p$};
				\end{tikzpicture}
		\end{center}
		\begin{itemchoice}
			\itemch Sai. Ta có \begin{align*}
			\overrightarrow{AB}\cdot \overrightarrow{CA}&=-\overrightarrow{CD}\cdot\overrightarrow{CA}\\ 
			&=-CD\cdot CA\cdot\cos\widehat{ACD}\\ 
			&=-a\sqrt{2}\cdot  a\cdot\cos45^\circ=-a^2.
		\end{align*}
		\itemch  Sai. Ta có $\overrightarrow{AM}\cdot \overrightarrow{AC}=AM\cdot AC\cos\widehat{MAC}=\dfrac{a}{2}\cdot a\sqrt{2}\cdot\cos 45^\circ=\dfrac{a^2}{2}$.
		\itemch Đúng. Ta có \begin{align*}
		 \overrightarrow{AD}\cdot \overrightarrow{BD}+\overrightarrow{O M}\cdot \overrightarrow{AC}&=\overrightarrow{DA}\cdot\overrightarrow{DB}-\dfrac{1}{2}\cdot\overrightarrow{AD}\cdot\overrightarrow{AC}\\
		&=a\cdot a\sqrt{2}\cos 45^\circ  -\dfrac{1}{2}\cdot a\cdot a\sqrt{2}\cdot\cos 45^\circ=\dfrac{a^2}{2}.
	\end{align*}
	\itemch Đúng. Ta có \begin{align*}
	 \left(\overrightarrow{AB}+\overrightarrow{AD}\right)\left(\overrightarrow{BD}+\overrightarrow{BC}\right)&=\overrightarrow{AC}\left(\overrightarrow{BD}+\overrightarrow{BC}\right)\\
	 &=\overrightarrow{AC}\cdot\overrightarrow{BD}+\overrightarrow{AC}\cdot\overrightarrow{BC}\\
	&=0+\overrightarrow{CA}\cdot\overrightarrow{CB}\\ 
	&=a\sqrt{2}\cdot  a\cdot\cos 45^\circ=a^2.
	\end{align*}
	\end{itemchoice}
} 
\end{ex}


%%%%%%%%%%%%%%%% TRẢ LỜI NGẮN %%%%%%%%%%%%%%%%

% nguồn: MaHoaNH_TraLoiNgan_Lop10
\begin{bt}%[0H5H4-1]%[ma:2fd9bf11]
	Cho tam giác $ABC$ vuông tại $A$ có cạnh $AC=7$ cm và $BC=14$ cm. Tính cô-sin của góc giữa hai vectơ $\overrightarrow{AC}$ và $\overrightarrow{CB}$. 
	\shortans{-0,5} 
	\loigiai{ 
	\immini{Ta có $\left(\overrightarrow{AC},  \overrightarrow{CB}\right)=180^{\circ}-\left(\overrightarrow{CA}, \overrightarrow{CB}\right)=180^{\circ}-\widehat{ACB}$.\\
		 Mà $\cos \widehat{ACB}=\dfrac{AC}{BC}=\dfrac{1}{2}$ nên $\widehat{ACB}=60^{\circ}$.\\
		 Do đó $\left(\overrightarrow{AC}, \overrightarrow{CB}\right)=180^{\circ}-60^{\circ}=120^{\circ}$.\\
		Vậy $\cos (\overrightarrow{AC}, \overrightarrow{CB})=\cos 120^{\circ}=-\dfrac{1}{2}=-0{,}5$.}{\begin{tikzpicture}[scale=1,font=\footnotesize,line join=round,line cap=round,>=stealth]
		\path (0,0) coordinate (A)++(1.5,0) coordinate  (C)
			 ($(C)!1!-60:(A)$) coordinate (c)
			(A)++(0,1)  coordinate (a)
			(intersection of A--a  and C--c) coordinate (B);
		 \draw (A)--(B)--(C)--cycle;
		\foreach \p/\g  in {A/210,B/90,C/-30} \fill[black] (\p) circle (1pt)+(\g:0.3) node{$\p$};
		 \draw pic[draw,angle radius=0.2cm]{right angle=B--A--C};
	\end{tikzpicture}
	 }}
\end{bt}

% nguồn: MaHoaNH_TraLoiNgan_Lop10
\begin{bt}%[0H5H4-1]%[ma:34671b90]
	Cho tam giác $ABC$ vuông tại $A$, trên hai cạnh  $AB$ và $AC$ lần lượt lấy hai điểm $B'$ và $C'$ sao cho $AB \cdot AB'=AC \cdot AC'$. Gọi $M$ là trung điểm của $BC$. Tính $\overrightarrow{AM}\cdot \overrightarrow{B'C'}$
	\shortans{0} 
	 \loigiai{
	\immini{Vì $M$ là  trung điểm của $BC$ nên $\overrightarrow{AM}=\dfrac{1}{2}\left(\overrightarrow{AB}+\overrightarrow{AC}\right)$.\\
		Do đó, \begin{align*}
		2\overrightarrow{AM}\cdot\overrightarrow{B'C'}&=\left(\overrightarrow{AB}+\overrightarrow{AC}\right)\left(\overrightarrow{AC'}-\overrightarrow{AB'}\right)\\ 
		 &=\overrightarrow{AB}\cdot\overrightarrow{AC'}+\overrightarrow{AC}\cdot\overrightarrow{AC'}-\overrightarrow{AB}\cdot\overrightarrow{AB'}-\overrightarrow{AC}\cdot\overrightarrow{AB'}\\
		&=0-AB\cdot  AB'+AC\cdot AC'=0.
	\end{align*} }{\begin{tikzpicture}[scale=1,font=\footnotesize,line join=round,line cap=round,>=stealth]
		\path (0,0)  coordinate (A)++(3,0) coordinate (C)
			 (A)++(0,2) coordinate (B)
			($(B)!0.5!(C)$)  coordinate (M);
		 \draw (A)--(B)--(C)--cycle;
		\foreach \p/\g in {A/210,B/90,C/-30,M/45} \fill[black] (\p) circle  (1pt)+(\g:0.3) node{$\p$};
		\draw pic[draw,angle  radius=0.2cm]{right angle=B--A--C};
	\end{tikzpicture}
	} 
} 
\end{bt}

% nguồn: MaHoaNH_TraLoiNgan_Lop10
\begin{bt}%[0H5H4-1]%[ma:e148cfa6]
	Cho tam giác $ABC$ vuông tại $A$ có $AB=a, AC=2\sqrt 3 a$ và $AM$ là trung tuyến. Tính tích vô hướng $\overrightarrow{BA} \cdot  \overrightarrow{AM}$ có dạng $\dfrac{m}{n}a^2,\, \dfrac{m}{n}$ là phân số tối giản và $m \in \mathbb{Z}, n \in \mathbb{N}$. Khi đó $m^2+n^2$ bằng bao nhiêu?
	\begin{center}
		\begin{tikzpicture}[scale=1,font=\footnotesize,line join=round,line cap=round,>=stealth]
				\path 
					(-2,0) coordinate (C) 
					(2,0) coordinate (B) 
					(60:2) coordinate (A) 
					 ($(B)!0.5!(C)$) coordinate (M);;
				\draw (A)--(C)--(B)--(A)--(M); 
				\draw pic[draw,angle radius = 6pt]{right angle = C--A--B}; 
				\foreach \p/\q in {A/90,C/-90,B/-90,M/-90} \fill[black] (\p)  circle (1pt)
					($(\p)+(\q:2.5mm)$)  node{$\p$};
			\end{tikzpicture}
	\end{center}
	\shortans{5} 
	 \loigiai{
		Ta  có $BC=\sqrt{AB^2+AC^2}=2a$, suy ra $MB=MC=MA=a$.\\
		Tam giác  $AMB$ có $AM=BM=AB$ nên là tam giác $AMB$ đều. Suy ra $\widehat{MAB}=60^\circ$.\\
		Khi đó $\overrightarrow{BA} \cdot \overrightarrow{AM}=-\overrightarrow{AB} \cdot \overrightarrow{AM}=-\left|\overrightarrow{AB}\right| \cdot\left|\overrightarrow{AM}\right| \cdot \cos \left(\overrightarrow{AB}, \overrightarrow{AM}\right)=-a \cdot a \cdot \cos 60^\circ=-\dfrac{a^2} 2$.\\ 
		Vậy  $m^2+n^2=\left(-1\right)^2+2^2=5$.
	} 
\end{bt}

% nguồn: MaHoaNH_TraLoiNgan_Lop10
\begin{bt}%[0H5H4-1]%[ma:e450de58]
	Cho  tam giác $ABC$ cân tại $A$, $M$ là trung điểm của $BC$, $H$ là hình chiếu của $M$ trên $AC$, $E$ là trung điểm của $MH$. Tính $\overrightarrow{AE} \cdot \overrightarrow{BH}$.
	\shortans{0} 
	 \loigiai{
	\immini{ 
		Ta có biến đổi tích vô  hướng như sau
		\begin{eqnarray*}
			2 \overrightarrow{AE} \cdot \overrightarrow{BH}&=&(\overrightarrow{AM}+\overrightarrow{AH}) \cdot\left(\overrightarrow{BM}+\overrightarrow{MH}\right)\\ 
			&=&\overrightarrow{AM} \cdot  \overrightarrow{MH}+\overrightarrow{AH} \cdot \overrightarrow{BM}\\
			&=&\overrightarrow{AM} \cdot \overrightarrow{MH}+\left(\overrightarrow{AM}+\overrightarrow{MH}\right) \cdot \overrightarrow{BM}\\ 
			&=&\overrightarrow{AM} \cdot \overrightarrow{MH}+\overrightarrow{MH} \cdot \overrightarrow{MC}\\ 
			&=&\overrightarrow{HM} \cdot \overrightarrow{MH}+\overrightarrow{MH}  \cdot \overrightarrow{MH}\\
			 &=&-\overrightarrow{MH}^2+\overrightarrow{MH}^2=0.
		\end{eqnarray*}
		Suy ra $\overrightarrow{AE}  \cdot \overrightarrow{BH}=0$. }{
		\begin{tikzpicture}[line join = round,line cap = round,>=stealth,font=\footnotesize,scale=.9]
				\coordinate (B) at ($(0,0)$); 
				\coordinate (C)  at ($(B)+(4,0)$);
				 \coordinate (M) at ($(B)!0.5!(C)$);
				\coordinate (A) at ($(M)+(0,4)$); 
				\coordinate (H) at  ($(C)!(M)!(A)$);
				 \coordinate (E) at ($(M)!0.5!(H)$);
				\draw (A)--(B)--(C)--(A)--(M) 
					(B)--(H)--(M) 
					 (A)--(E);
				\foreach  \x/\y in {A/180,B/-180,C/30,M/-90,H/0,E/-90}
				\draw[fill] (\x) circle (1pt) + (\y:0.25cm) node{$\x$}; 
				\path  pic[draw,angle radius=6]{right angle=M--H--C};
				\path pic[draw,angle radius=6]{right angle=A--M--B}; 
			\end{tikzpicture}
		} 
	} 
\end{bt}

% nguồn: MaHoaNH_TraLoiNgan_Lop10
\begin{bt}%[0H5H4-1]%[ma:656b9dc7]
	Cho hình thang vuông $ABCD$ có đáy lớn $AB=8a$, đáy nhỏ $CD=4a$, đường cao $AD=6a$,  $I$ là trung điểm của $AD$. Biết giá trị $\left(\overrightarrow{IA}+\overrightarrow{IB}\right) \cdot \overrightarrow{ID}$ có dạng $ka^2$. Khi đó $k$ bằng bao nhiêu?
	\shortans{-18} 
	\loigiai{ 
	 \immini{
		\begin{eqnarray*}
			\left(\overrightarrow{IA}+\overrightarrow{IB}\right) \cdot  \overrightarrow{ID}
			&=&  \overrightarrow{IA}\cdot \overrightarrow{ID}+\overrightarrow{IB}\cdot \overrightarrow{ID} \\
			 &=&-{\overrightarrow{IA}}^2+IB\cdot ID\cdot \cos \widehat{BID} \\
			&=&-IA^2-IB\cdot ID\cdot \cos \widehat{BIA} \\ 
			 &=&-IA^2-IB\cdot ID\cdot \dfrac{IA}{IB} \\
			 &=&-IA^2-IA^2\\
			 &=&-2IA^2=-18a^2.
		\end{eqnarray*}
		 }{
		\begin{tikzpicture}[scale=0.6,font=\footnotesize,line join=round,line cap=round,>=stealth]
				\def\a{1} 
				 \coordinate[label=below:$A$] (A) at (0,0);
				\coordinate[label=above:$D$]  (D) at (0,6*\a);
				\coordinate[label=below:$B$] (B)  at (8*\a,0);
				\coordinate[label=above:$C$] (C)  at (4*\a,6*\a);
				 \coordinate[label=left:$I$] (I) at ($(A)!0.5!(D)$);
				\node  at (4*\a,0) [below]{$8a$};
				\node at (-0.5*\a,3.5*\a) [above left]{$6a$}; 
				 \node at (2*\a,6*\a) [below]{$4a$};
				 \draw[->,arrows={-latex}] (I)--(A);
				\draw[->,arrows={-latex}] (I)--(B); 
				 \draw[->,arrows={-latex}] (I)--(D);
				 \draw (A)--(B)--(C)--(D)--cycle;
				\foreach \diem in {A,B,C,D,I}  \fill (\diem) circle (1.2pt);
			\end{tikzpicture}
		 }
	 }
\end{bt}


%%%%%%%%%%%%%%%% TỰ LUẬN %%%%%%%%%%%%%%%%

% nguồn: MaHoaNH_TuLuan_Lop10
\begin{bt}%[0H5H4-1]%[ma:98f4864f]
	Cho hình thoi $ABCD$ tâm $O$ có  cạnh bằng $4$ và $\widehat{ABD}=60^{\circ}$. Gọi $I$ là điểm thỏa mãn $2 \overrightarrow{IC}+\overrightarrow{ID}=\overrightarrow{0}$. Tính tích vô hướng $\overrightarrow{AO}\cdot \overrightarrow{BI}$.
	\loigiai{ 
		 Hình vẽ
		\begin{center}
			\begin{tikzpicture}[scale=1,font=\footnotesize,line join=round,line cap=round,>=stealth]
					\path (0,0) coordinate (A)++(4,0) coordinate  (C)
						 ($(A)!0.5!(C)$) coordinate (O)
						($(A)!1!30:(C)$) coordinate  (a)
						(O)++(0,1) coordinate (o) 
						(intersection of A--a and  O--o) coordinate (B)
						($(A)+(C)-(B)$) coordinate (D) 
						($(C)!1/3!(D)$) coordinate (I); 
					\draw (A)--(B)--(C)--(D)--cycle  (A)--(C)
						 (I)--(B)--(D);
					\foreach \p/\g in {A/180,B/90,C/0,D/-90,I/-45,O/225} \fill[black]  (\p) circle (1pt)+(\g:0.3) node{$\p$};
				\end{tikzpicture}
		\end{center}
		Do $ABCD$ là hình thoi có cạnh bằng $4$ và  $\widehat{ABD}=60^{\circ}$ nên $ABD$ và $BCD$ là các tam giác đều cạnh $4$.\\
		Ta có \begin{align*}
		\overrightarrow{AO}\cdot\overrightarrow{BI}&=\overrightarrow{AO}\cdot(\overrightarrow{BD}+\overrightarrow{DI})=\overrightarrow{AO}\cdot\overrightarrow{DI}\\ 
		 &=\overrightarrow{AO}\cdot\left(\dfrac{2}{3}\overrightarrow{DC}\right)=\dfrac{2}{3}\overrightarrow{AO}\cdot\overrightarrow{AB}\\
		 &=\dfrac{2}{3}\cdot\dfrac{4\sqrt{3}}{2}\cdot 4\cdot\cos 30^{\circ}=8.
	\end{align*}
 }
\end{bt}
